Offloading tasks to edge servers or the cloud trades completion time against device energy and cost, under uncertain
bandwidth and server capacity. We formulate the offloading of directed acyclic graph (DAG) workflows over mobile devices,
edge servers and a cloud as an interval multi-objective problem: task placement and task order are optimised for
makespan along the interval critical path, device energy and cost, with uncertain quantities given as intervals.
Extending a recent reinforcement-learning-driven interval evolutionary algorithm to this problem, we show that its
convergence ranking, an interval dominance degree, stops selecting once most of the population is mutually
non-dominated, and that its learned control of the selection weights does not beat a fixed choice because the
per-generation reward is almost pure noise. The proposed \ours{} uses an interval best-position Tchebycheff rank as the
convergence tier of the dominance-, diversity- and uncertainty-aware three-tier selection, mates neighbours in
objective space and uses DAG-valid crossover and mutation. On 57 instances with up to 1000 tasks, it is compared with
four interval evolutionary algorithms, NSGA-II, NSGA-III, the reinforcement-learning-driven predecessor and
list-scheduling heuristics. \ours{} has the best IGD and hypervolume on all 12 instances of the predecessor's sizes and
beats the strongest baseline on 43 of 45 workflow instances in IGD and 42 in hypervolume, at a similar running time.
At 1000 and 2000 tasks it keeps the higher hypervolume, while IGD is level. A Monte Carlo check confirms that the
predicted objective intervals contain the realised values.
